In this section, we investigate obsolescence and analyze the lifespan structure of individual highly cited papers to assess whether they sustain long-term influence. To identify papers that have reached an obsolescent state, we draw upon the long-term citation model of P. D. B. Parolo et al. \cite{Par1} and the citation decay model of D. Wang et al. \cite{Wan1}, and define the following three criteria: (1) a sufficiently long time has elapsed since publication, (2) the citation peak has passed by a sufficient margin, and (3) citation activity has decayed sufficiently. The specific threshold values used for these criteria are presented in Figure \ref{fig:gr6}.
